\begin{table}[htbp]
\centering
\footnotesize
\setlength{\tabcolsep}{3.5pt}
\renewcommand{\arraystretch}{0.90}
\caption{Trade Union Organization, Density over Employed Workers, and Average Union Size (1932--1975)}
\label{tab:union_density_osorio_polanco}
\begin{tabularx}{\textwidth}{lcccccX}
\toprule
\textbf{Year} & \textbf{Unions} & \textbf{Affiliated} & \textbf{Employed} & \textbf{Avg Size} & \textbf{Density (\%)} & \textbf{Accumulation Regime \& Institutional Era} \\
\midrule
1932 & 421 & 54,801 & 1,165,000 & 130.2 & 4.70\% & Post-Depression trough; Alessandri restoration \\
1935 & 669 & 83,262 & 1,553,000 & 124.5 & 5.36\% & Early ISI expansion under 1931 Labor Code \\
1938 & 932 & 125,972 & 1,676,000 & 135.2 & 7.52\% & Popular Front victory (Aguirre Cerda) \\
1939 & 1,687 & 173,438 & 1,709,000 & 102.8 & 10.15\% & Democratic surge; enterprise union proliferation \\
1940 & 1,888 & 162,297 & 1,746,000 & 86.0 & 9.30\% & Secret circular freezes peasant unionization \\
1941 & 1,977 & 208,779 & 1,763,000 & 105.6 & 11.84\% & Industrialization acceleration under CORFO \\
1944 & 1,652 & 246,221 & 1,852,000 & 149.0 & 13.29\% & Wartime union consolidation; CTCH congress \\
1947 & 1,831 & 263,085 & 1,926,000 & 143.7 & \textbf{13.66\%} & Pre-\textit{Ley Maldita} peak of labor mobilization \\
1948 & 1,857 & 263,676 & 2,029,000 & 142.0 & \textbf{13.00\%} & Enactment of \textit{Ley Maldita} (Laws 8,987 \& 8,811) \\
1950 & 1,907 & 260,143 & 2,087,000 & 136.4 & 12.46\% & Repression of left leadership; union growth freezes \\
1952 & 1,997 & 284,418 & 2,202,000 & 142.4 & 12.92\% & Ibáñez del Campo populist election \\
1955 & 2,177 & 305,192 & 2,209,000 & 140.2 & 13.82\% & Klein-Saks mission wage freeze \\
1958 & 1,894 & 276,346 & 2,297,000 & 145.9 & 12.03\% & Repeal of \textit{Ley Maldita}; density stagnation \\
1960 & 1,770 & 272,966 & 2,313,000 & 154.2 & 11.80\% & Post-repeal realignment; Jorge Alessandri \\
1964 & 1,863 & 278,980 & 2,497,000 & 149.7 & \textbf{11.17\%} & Frei Montalva election; rural baseline \\
1966 & 2,882 & 369,507 & 2,604,000 & 128.2 & 14.19\% & Peasant organizing begins (\textit{Promoción Popular}) \\
1967 & 3,426 & 442,650 & 2,681,000 & 129.2 & 16.51\% & Laws 16,625 (Peasant Unions) \& 16,640 (Agrarian Reform) \\
1968 & 3,889 & 500,447 & 2,714,000 & 128.7 & 18.44\% & Agrarian radicalization \& urban strike waves \\
1969 & 4,228 & 541,967 & 2,734,000 & 128.2 & 19.82\% & Pre-electoral polarization; \textit{El Tacnazo} \\
1970 & 4,581 & 627,664 & 2,743,000 & 137.0 & \textbf{22.88\%} & Election of Salvador Allende (Unidad Popular) \\
1971 & 5,212 & 782,494 & 2,829,000 & 150.1 & 27.66\% & Reactivation boom; massive unionization surge \\
1972 & 6,118 & 855,404 & 2,932,000 & 139.8 & \textbf{29.17\%} & Distributive peak; Bosses' Lockout (\textit{Paro de Octubre}) \\
1973 & 6,692 & 939,319 & 2,976,000 & 140.4 & \textbf{31.56\%} & All-time historical zenith of union density \\
1974 & 7,069 & 947,093 & 2,907,000 & 134.0 & 32.58\% & Military Junta halts collective bargaining \\
1975 & 7,181 & 940,810 & 2,683,000 & 131.0 & 35.07\% & Shock therapy depression; employment contraction \\
\bottomrule
\end{tabularx}
\begin{minipage}{\textwidth}
\vspace{0.3em}
\footnotesize
\textit{Source:} \citet{OsorioPolanco2021}, \textit{Repositorio de Estadísticas Sindicales (RES)}, sheet \texttt{Empalme}. Average union size is $\text{Afiliación} / \text{Sindicatos}$. Union density is $\text{Afiliación} / \text{Ocupados}$.
\end{minipage}
\end{table}
